\clearpage\fancyhead[L]{\small GAUSS TO VON LINDENAU.}\SourceLeaf{330}
\begin{center}\textbf{GAUSS TO VON LINDENAU.}\par\vspace{4pt}\hrule\vspace{4pt}\fontsize{10}{12}\selectfont Monatliche Correspondenz. Vol. XXIII. pp. 97. 98. January 1811.\par\vspace{4pt}\hrule\vspace{4pt}\end{center}

\begin{flushright}Göttingen, 14 January 1811.\end{flushright}
\par\indent --- --- --- At this observatory, the night of 4 to 5 January was the first in nearly two months on which it was possible to try to find \emph{Pallas}. As I could use only movable instruments, the moonlight and severe cold (11$^\circ$ Réaum.) made the search extremely difficult. Even so, my efforts were rewarded. I found \emph{Pallas} where it had to be according to my corrected ephemeris and also managed to make a couple of comparisons with stars from the \emph{Histoire céleste}. From these I derived the following position of the planet, subject to a more precise reduction later.

\begin{center}\fontsize{10}{12}\selectfont\setlength{\tabcolsep}{4pt}\renewcommand{\arraystretch}{1.15}\begin{tabular}{c|c|c|c}
1811 & Mean time & \shortstack{Right asc.\\of Pallas} & \shortstack{Southern\\declination}\\\hline
4 January & $15^{\text{u}}\;4^{\text{m}}\;42^{\text{s}}$ & $151^\circ\;40'\;5''$ & $22^\circ\;22'\;36''$\\
\end{tabular}\end{center}

\Needspace{16\baselineskip}\par\indent My newly revised calculations of the perturbations have given me the following elements of \emph{Pallas}. They differ only slightly from those I last communicated to you (Dec. issue 1810, p. 592) [p. 321 of this volume].

\begin{center}\fontsize{10}{12}\selectfont\setlength{\tabcolsep}{4pt}\renewcommand{\arraystretch}{1.15}\begin{tabular}{llr}
\shortstack[l]{Epoch of mean longitude\\on the Göttingen meridian} & .\;.\;. & $126^\circ\;33'\;6''$\\
Longitude of perihelion & .\;.\;.\;.\;.\;.\;.\;.\;.\;.\;.\;.\;.\;.\;. & $120\;55\;8$\\
Longitude of the ascending node & .\;.\;.\;.\;.\;.\;.\;.\;.\;.\;. & $172\;32\;40$\\
Inclination of the orbit & .\;.\;.\;.\;.\;.\;.\;.\;.\;.\;.\;.\;.\;.\;.\;. & $34\;35\;14$\\
Mean daily tropical motion & .\;.\;.\;.\;.\;.\;.\;.\;. & $768.980$\\
Eccentricity & .\;.\;.\;.\;.\;.\;.\;.\;.\;.\;.\;.\;.\;.\;.\;.\;.\;. & $0.24160$\\
Logarithm of the semimajor axis & .\;.\;.\;.\;.\;.\;.\;.\;. & $0.44278$\\
\end{tabular}\end{center}

\par\indent Here is the calculated continuation of the ephemeris [from 3 Feb. to 12 April 1811].
